% ---
% titre: Aire d'un losange à partir du périmètre
% theme: GM
% chapitre: Lignes et surfaces
% sous_chapitre: Lignes et surfaces
% niveau: [9e]
% type: EVAL
% tags: [c-perimetre, c-aire, c-losange, c-parallelogramme, p-question-TS]
% difficulte: 2
% corrige: true
% ---

\question[5]
Le périmètre du losange ci-dessous est égal à 168 mm. Calcule son aire.

\vspace{10pt}
\begin{center}
\begin{tikzpicture}[scale=0.8, rotate=20]
	\tkzDefPoints{0/0/A,1.5/4/B,5.5/4/C,4/0/D,2/0/H,2/4/I}
	\tkzDrawPolygon(A,B,C,D)
	\tkzDrawSegment[dim={{3,5\,cm},-20pt,midway}, dim style/.append style={black,dashed},dim fence style/.append style={draw=none}](H,I)
	\tkzMarkRightAngle[size=0.5,draw](A,H,I)
\end{tikzpicture}
\end{center}

\vspace{10pt}{\footnotesize\raggedleft Espace pour ta démarche et tes calculs\par}
\begin{solutionorgrid}[150pt]
Côté du losange : $168 \div 4 = 42\,mm = 4,2\,cm$ \hfill \textbf{(2pt)}

\vspace{10pt}
Un losange est un parallélogramme particulier : son aire se calcule donc avec $A = c\cdot h$, où $c$ est un côté et $h$ la hauteur correspondante. \hfill

\vspace{10pt}
$A =  c\cdot h = 4,2\cdot 3,5 = 14,7\,cm^2$ \hfill \textbf{(2,5pt)}
\end{solutionorgrid}

\begin{tcolorbox}[enhanced jigsaw, interior hidden, colframe=box, parbox=false]%
Réponse: 
\sol{\vspace{20pt}}{L'aire du losange vaut 14,7\,cm$^2$. \textbf{(0,5pt)}}
\end{tcolorbox}